\phantomsection\addcontentsline{toc}{section}{High-Speed Transportation}
\ECchapter{30}{High-Speed Transportation}{Moving faster with less energy}{ECBlue}

\ECObjectives{
\begin{itemize}
\item Explain how drag, power and braking distance scale with speed.
\item Compare conventional rail, magnetic levitation and low-pressure concepts.
\item Evaluate travel time, energy use and system-level safety constraints.
\end{itemize}}

\begin{multicols}{2}
\begin{engineeringnews}
High-speed rail, magnetic-levitation systems and low-pressure-tube concepts seek to reduce journey times between cities. At high speed, aerodynamic drag, infrastructure quality, signalling, energy supply and emergency procedures become dominant engineering constraints.
\end{engineeringnews}

\begin{ecarticle}
High-speed transportation is a system problem rather than simply a faster vehicle. Track geometry, propulsion, signalling, stations, maintenance and emergency procedures must be designed together. A train travelling at 300 km/h moves at $83.3$ m/s, so it covers the length of a football pitch in little more than one second.

Aerodynamic drag can be approximated by
\[
F_D=\frac{1}{2}\rho C_D A v^2,
\]
where $\rho$ is air density, $C_D$ the drag coefficient, $A$ the frontal area and $v$ the speed. The corresponding power is
\[
P_D=F_Dv=\frac{1}{2}\rho C_D A v^3.
\]
Therefore, doubling speed multiplies aerodynamic power by about eight if geometry and air density remain unchanged. Streamlined noses, smooth inter-car gaps and protected underbodies reduce drag and pressure waves. Tunnel entry is particularly demanding because the train compresses air in a confined space.

Conventional high-speed trains use steel wheels on rails. This gives low rolling resistance, mature switches and compatibility with parts of existing networks. Magnetic levitation removes wheel--rail contact at speed and can reduce mechanical wear, but it requires dedicated guideways, accurate gap control and specialised power electronics. Low-pressure tubes could reduce aerodynamic resistance further, yet maintaining long sealed infrastructure, controlling thermal expansion and providing safe evacuation remain major challenges.

For a constant deceleration $a$, the ideal braking distance is
\[
d=\frac{v^2}{2a}.
\]
Real stopping distance is longer because it includes detection, command and brake-build-up delays. High-speed networks therefore use continuous cab signalling and automatic train protection rather than relying on lineside signals. Regenerative braking can return energy to the supply system when another train, the grid or a storage device can absorb it.

Maximum speed does not determine usefulness on its own. Door-to-door time includes access, waiting, intermediate stops and transfers. Capacity, reliability, construction carbon, noise and ticket cost also matter. High-speed rail performs best on busy corridors where high passenger flows justify dedicated infrastructure. Engineers should compare energy per passenger-kilometre and lifecycle impacts, not only vehicle speed or peak power.
\end{ecarticle}

\begin{tcolorbox}[title=\sffamily\bfseries Did You Know?,colback=yellow!10,colframe=ECOrange]
Increasing maximum speed may save only a few minutes when stations are close together, because acceleration, braking and dwell times occupy a large fraction of the journey.
\end{tcolorbox}

\begin{engineeringfocus}
\textbf{Energy conversion in high-speed rail}
\begin{center}
\resizebox{0.96\linewidth}{!}{%
\begin{tikzpicture}[font=\scriptsize,>=Latex,node distance=4mm]
\node[draw=ECBlue,fill=ECLightBlue,rounded corners] (grid) {electric grid};
\node[draw=ECPurple,fill=ECPurple!10,rounded corners,right=of grid] (drive) {converter and motor};
\node[draw=ECGreen,fill=ECGreen!10,rounded corners,right=of drive] (motion) {vehicle motion};
\node[draw=ECOrange,fill=ECOrange!10,rounded corners,below=of drive] (loss) {drag, rolling and auxiliaries};
\draw[->,thick] (grid)--(drive);
\draw[->,thick] (drive)--(motion);
\draw[->,thick] (motion)--(loss);
\draw[->,thick,dashed] (motion) to[bend left=35] node[above]{regeneration} (grid);
\end{tikzpicture}}
\end{center}
\end{engineeringfocus}

\begin{keyfigures}
\begin{tabularx}{\linewidth}{@{}Xr@{}}
Drag force & $\propto v^2$\\
Aerodynamic power & $\propto v^3$\\
Ideal braking distance & $v^2/(2a)$\\
Safety system & continuous signalling\\
Useful metric & passenger-km
\end{tabularx}
\end{keyfigures}

\begin{tcolorbox}[title=\sffamily\bfseries Aerodynamic power versus speed,colback=white,colframe=ECBlue]
\begin{center}
\begin{tikzpicture}
\begin{axis}[width=0.96\linewidth,height=46mm,xlabel={Speed (km/h)},ylabel={Relative aerodynamic power},grid=major,ticklabel style={font=\scriptsize},label style={font=\scriptsize}]
\addplot[thick,mark=*] table[x=speed,y=power,col sep=comma]{data/EC30/power.csv};
\end{axis}
\end{tikzpicture}
\end{center}
\end{tcolorbox}

\begin{tcolorbox}[title=\sffamily\bfseries Quantitative application,colback=ECBlue!5,colframe=ECBlue]
A train has $C_DA=8.0$ m$^2$ and runs in air of density $1.2$ kg/m$^3$.
\begin{enumerate}
\item Estimate aerodynamic drag and power at 300 km/h.
\item With an emergency deceleration of $0.9$ m/s$^2$, calculate the ideal braking distance from 300 km/h.
\item Add a 5 s response delay and estimate the total stopping distance.
\end{enumerate}
\textit{Expected result:} $F_D\approx33$ kN, $P_D\approx2.8$ MW, ideal braking distance $\approx3.86$ km and total distance $\approx4.28$ km.
\end{tcolorbox}

\begin{vocabulary}
\begin{tabularx}{\linewidth}{@{}lX@{}}
aerodynamic drag & traînée aérodynamique\\
rolling resistance & résistance au roulement\\
guideway & voie de guidage\\
magnetic levitation & sustentation magnétique\\
pressure wave & onde de pression\\
braking distance & distance de freinage\\
regenerative braking & freinage régénératif\\
signalling & signalisation\\
passenger-kilometre & passager-kilomètre\\
dedicated infrastructure & infrastructure dédiée
\end{tabularx}
\end{vocabulary}

\begin{questions}
\item Why does aerodynamic power rise more quickly than drag force?
\item What are the main advantages and limits of magnetic levitation?
\item Why are continuous signalling systems necessary?
\item Under what conditions can regenerative braking recover useful energy?
\item Why can maximum speed be a poor indicator of door-to-door performance?
\end{questions}

\begin{engineeringchallenge}
Compare a 300 km/h railway with a magnetic-levitation line for a 600 km corridor. Estimate non-stop travel time, identify dominant energy losses, define safety requirements and propose a lifecycle comparison including infrastructure carbon and passenger demand.
\end{engineeringchallenge}

\begin{speaking}
Should public investment favour maximum speed or a denser network with lower speeds and more frequent services?
\end{speaking}
\end{multicols}

\subsection*{Sources}
\ECsource{International Union of Railways -- High-speed rail}{https://uic.org/passenger/highspeed/}
\ECsource{International Energy Agency -- Rail}{https://www.iea.org/energy-system/transport/rail}
